\textbf{Context:} The growing adoption of Artificial Intelligence (AI) systems, particularly Large Language Models (LLMs), has intensified concerns regarding transparency, fairness, and accountability, especially due to documented incidents involving algorithmic bias, discriminatory automated decisions, and limited auditability. Although ethical guidelines and regulatory frameworks, such as those proposed by UNESCO, the OECD, and the EU AI Act, establish high-level principles, the literature indicates that their technical operationalization remains limited. In this context, a critical gap has been identified in the \textit{Testing} phase of the software lifecycle, characterized by a reliance on manual audits and a lack of automated tools capable of detecting ethical violations in deployed AI systems.
\textbf{Objective:} This dissertation aimed to develop and evaluate a set of automated \textit{fuzzing engines} capable of detecting ethical violations related to \textit{Fairness}, \textit{Accountability}, and \textit{Transparency} in deployed AI models, contributing to the development of tools that can be integrated into real-world Software Engineering workflows. The focus on testing is motivated by the need for model-agnostic, \textit{black-box} auditing mechanisms capable of validating the behavior of the final system independently of the training data used.
\textbf{Method:} The research followed a \textit{Design Science Research} (DSR) approach. First, a Systematic Mapping Study (2021--2025) was conducted to identify the most prevalent ethical principles and existing technical gaps. Subsequently, the eleven risks derived from the literature-based taxonomy were specified and mapped to testing strategies, of which six were operationalized into modular \textit{engines} employing five \textit{fuzzing} techniques: \textit{mutation-based}, \textit{generation-based}, \textit{differential}, \textit{metamorphic testing}, and adversarial prompt injection, in addition to deterministic oracles based on \textit{a priori} thresholds. Finally, the \textit{engines} were empirically evaluated through \textit{black-box} \textit{fuzzing} campaigns applied to three commercial LLMs (GPT, Gemini, and DeepSeek). To ensure reproducibility, the collection and evaluation phases were conducted separately, and failure rates were recorded by risk, provider, and principle.
\textbf{Results:} The mapping identified \textit{Fairness}, \textit{Accountability}, and \textit{Transparency} as the most frequently addressed principles in the literature and the ones exhibiting the highest degree of operational maturity. The results also highlighted the limited coverage of the \textit{Testing} phase (27.1\%) and the absence of automated tools suitable for integration into real-world environments. Of the eleven risks included in the taxonomy, six proved fully verifiable through automation, whereas five require complementary human evaluation due to their dependence on social context. The empirical evaluation, conducted on 8,880 pairs and approximately 14,160 API calls, revealed an overall ethical failure rate of 32.7\% (95\% CI [31.7; 33.7]), with the following results by model: GPT (30.3\%), Gemini (32.7\%), and DeepSeek (35.1\%). Failures were concentrated in \textit{Fairness}, particularly in counterfactual consistency (R-F1: 79.2\%), and in sensitivity to irrelevant attributes (R-T2: 70.0\%). When aggregated by principle, failure rates were highest for \textit{Transparency} (39.9\%), followed by \textit{Accountability} (30.8\%) and \textit{Fairness} (30.1\%). In explanatory \textit{Transparency} (R-T1), a pronounced asymmetry was observed between direct explanation (22--27\% failure rate) and meta-explanation (0--1.5\%).
\textbf{Conclusion:} The resulting artifacts reduce the gap between abstract ethical guidelines and practical verification mechanisms, enabling the systematic detection of ethical failures in AI models. Furthermore, they eliminate the need for subjective evaluators, such as human inspectors or LLMs-as-a-judge, by relying on deterministic and reproducible oracles. This work also advances the state of the art by proposing an automated and modular approach to ethical testing, providing integrated coverage of the three FAT principles and directly addressing the gaps identified in the systematic mapping study.